\scanpage{29}
\mspage{3}
déterminée par la condition que l'aire totale de $X$ soit $=1$.

On trouve une \uline{métrique plate} caractérisée par la condition d'être invariante par translation, et d'aire \deleted{\unsure{volume}} totale égale à $1$ - donc stable par tt autom. conforme (pas seulement holomorphe). Je \unsure{sousent} que dans le cas non orienté, la métrique canonique de $X^{or}$ (qui est une courbe elliptique / $\mathbb C$) se descend à $X^{\unsure{or|an}}$.

$2'$) Cas $X\simeq\mathbb C^*$, ou\inserted{$X\simeq$} plan projectif réel privé d'un point [i.e. $X$ non orientable et $X^{or}\simeq\mathbb C^*$]. Comme dans $2$, dans le cas $X\simeq\mathbb C^*$, \deleted{plus généralement $X^{\unsure{or}}$} \deleted{\unsure{par}} descente à partir du rev. universel, on trouve une métrique canonique mod. multiplication par un scalaire $>0$, \uline{métrique plate} également.

\deleted{Comme} les autom. de $\mathbb C^*$\inserted{sont} \deleted{\unsure{précisément}}\inserted{\unsure{tt}} de la forme $z\longrightarrow az$, soit $z\longrightarrow a/z$, qui conservent la métrique [car sur $\widetilde X\simeq\mathbb C$ ils correspondent à des autom. $z\longrightarrow z+\alpha$, $z\longmapsto-z+\alpha$], Donc en \unsure{définissant|choisissant} la métrique sur $\mathbb C^*$, \deleted{\unsure{de façon}}\inserted{on}
% archive foliation: 31
